\documentclass[11pt]{article}
\usepackage[margin=1.1in]{geometry}
\usepackage{amsmath,amssymb,amsthm,mathtools,booktabs,array,tabularx,longtable}
\usepackage[dvipsnames]{xcolor}
\usepackage[colorlinks=true,linkcolor=MidnightBlue,citecolor=MidnightBlue,urlcolor=MidnightBlue]{hyperref}
\usepackage{fancyhdr}
\pagestyle{fancy}\fancyhf{}\renewcommand{\headrulewidth}{0pt}
\fancyfoot[R]{\thepage}
\fancyfoot[L]{\scriptsize\color{gray}Generated by the Transport Lab from approved parts of its store; not a deposited paper. Papers of record: \href{https://zenodo.org/communities/transport-papers}{Zenodo, Transport Papers}}
\fancypagestyle{plain}{\fancyhf{}\fancyfoot[R]{\thepage}\fancyfoot[L]{\scriptsize\color{gray}Generated by the Transport Lab from approved parts of its store; not a deposited paper. Papers of record: \href{https://zenodo.org/communities/transport-papers}{Zenodo, Transport Papers}}\renewcommand{\headrulewidth}{0pt}}
\newtheorem{theorem}{Theorem}[section]
\newtheorem{proposition}[theorem]{Proposition}
\newtheorem{lemma}[theorem]{Lemma}
\newtheorem{corollary}[theorem]{Corollary}
\newtheorem{conjecture}[theorem]{Conjecture}
\newtheorem{refutation}[theorem]{Refutation}
\newtheorem{prediction}[theorem]{Prediction}
\newtheorem{principle}[theorem]{Principle}
\newtheorem{observation}[theorem]{Observation}
\newtheorem{axiom}[theorem]{Axiom}
\theoremstyle{definition}
\newtheorem{definition}[theorem]{Definition}
\newtheorem{fact}[theorem]{Fact}
\newtheorem{assembled}[theorem]{Assembly}
\newtheorem{postulate}[theorem]{Postulate}
\newtheorem{premise}[theorem]{Premise}
\newtheorem{hypothesis}[theorem]{Hypothesis}
\newtheorem{remark}[theorem]{Remark}
\newtheorem{example}[theorem]{Example}
\theoremstyle{plain}
\setlength{\marginparwidth}{0.8in}
\newcommand{\changedmark}[2]{\marginpar{\raggedright\scriptsize\color{#1}#2}}
\providecommand{\Aa}{}\renewcommand{\Aa}{\mathcal{A}}
\providecommand{\Ad}{}\renewcommand{\Ad}{\widetilde{\mathrm{Ad}}}
\providecommand{\Aev}{}\renewcommand{\Aev}{\mathcal{A}^{+}}
\providecommand{\Aod}{}\renewcommand{\Aod}{\mathcal{A}^{-}}
\providecommand{\B}{}\renewcommand{\B}{\mathrm{B}}
\providecommand{\C}{}\renewcommand{\C}{\mathbb{C}}
\providecommand{\CAR}{}\renewcommand{\CAR}{\mathrm{CAR}}
\providecommand{\CS}{}\renewcommand{\CS}{{\mathbb{S}^{+}}}
\providecommand{\CSm}{}\renewcommand{\CSm}{{\mathbb{S}^{-}}}
\providecommand{\Cl}{}\renewcommand{\Cl}{\mathrm{Cl}}
\providecommand{\Con}{}\renewcommand{\Con}{\mathcal{C}}
\providecommand{\Dl}{}\renewcommand{\Dl}{\Delta}
\providecommand{\EN}{}\renewcommand{\EN}{E_{N}}
\providecommand{\End}{}\renewcommand{\End}{\mathrm{End}}
\providecommand{\F}{}\renewcommand{\F}{\mathbb{F}}
\providecommand{\G}{}\renewcommand{\G}{\Gamma}
\providecommand{\HH}{}\renewcommand{\HH}{\mathbb{H}}
\providecommand{\IS}{}\renewcommand{\IS}{I\sigma_3}
\providecommand{\J}{}\renewcommand{\J}{\mathbf{J}}
\providecommand{\LH}{}\renewcommand{\LH}{L_{H}}
\providecommand{\Mat}{}\renewcommand{\Mat}{\mathrm{M}}
\providecommand{\Mo}{}\renewcommand{\Mo}{M^{\circ}}
\providecommand{\Mp}{}\renewcommand{\Mp}{M^{+}}
\providecommand{\Msun}{}\renewcommand{\Msun}{M_{\odot}}
\providecommand{\Mt}{}\renewcommand{\Mt}{\widetilde{M}}
\providecommand{\Ncone}{}\renewcommand{\Ncone}{\mathcal{N}}
\providecommand{\Neff}{}\renewcommand{\Neff}{N_{\mathrm{eff}}}
\providecommand{\PCS}{}\renewcommand{\PCS}{\mathbb{P}{\mathbb{S}^{+}}}
\providecommand{\PT}{}\renewcommand{\PT}{\mathbb{PT}}
\providecommand{\Pin}{}\renewcommand{\Pin}{\mathrm{Pin}}
\providecommand{\Q}{}\renewcommand{\Q}{\mathbb{Q}}
\providecommand{\R}{}\renewcommand{\R}{\mathbb{R}}
\providecommand{\Ric}{}\renewcommand{\Ric}{\mathrm{Ric}}
\providecommand{\SBH}{}\renewcommand{\SBH}{S_{\mathrm{BH}}}
\providecommand{\Sc}{}\renewcommand{\Sc}{\mathbb{S}}
\providecommand{\Sch}{}\renewcommand{\Sch}{\mathrm{P}}
\providecommand{\Sig}{}\renewcommand{\Sig}{\Sigma}
\providecommand{\Som}{}\renewcommand{\Som}{S_\omega}
\providecommand{\Sp}{}\renewcommand{\Sp}{\mathbb{S}}
\providecommand{\Sph}{}\renewcommand{\Sph}{\mathrm{S}}
\providecommand{\Spi}{}\renewcommand{\Spi}{S_\pi}
\providecommand{\Spin}{}\renewcommand{\Spin}{\mathrm{Spin}}
\providecommand{\Sth}{}\renewcommand{\Sth}{S_{\mathrm{th}}}
\providecommand{\T}{}\renewcommand{\T}{\mathbb{T}}
\providecommand{\TH}{}\renewcommand{\TH}{T_{\mathrm{H}}}
\providecommand{\Tr}{}\renewcommand{\Tr}{\operatorname{Tr}}
\providecommand{\UHF}{}\renewcommand{\UHF}{\mathrm{UHF}}
\providecommand{\Z}{}\renewcommand{\Z}{\mathbb{Z}}
\providecommand{\Zc}{}\renewcommand{\Zc}{\mathrm{Z}}
\providecommand{\cmg}{}\renewcommand{\cmg}{\mathrm{cm^{2}\,g^{-1}}}
\providecommand{\dS}{}\renewcommand{\dS}{\mathrm{dS}}
\providecommand{\dsc}{}\renewcommand{\dsc}{\mathrm{dsc}}
\providecommand{\dtPG}{}\renewcommand{\dtPG}{dt}
\providecommand{\dtS}{}\renewcommand{\dtS}{dt_{\mathrm{S}}}
\providecommand{\ellv}{}\renewcommand{\ellv}{\ell}
\providecommand{\er}{}\renewcommand{\er}{e_r}
\providecommand{\et}{}\renewcommand{\et}{e_t}
\providecommand{\ev}{}\renewcommand{\ev}{\mathrm{ev}}
\providecommand{\gl}{}\renewcommand{\gl}{\mathfrak{gl}}
\providecommand{\hatV}{}\renewcommand{\hatV}{\widehat{V}}
\providecommand{\hb}{}\renewcommand{\hb}{\bar{h}}
\providecommand{\hf}{}\renewcommand{\hf}{\hat{f}}
\providecommand{\id}{}\renewcommand{\id}{\mathrm{id}}
\providecommand{\kpc}{}\renewcommand{\kpc}{\mathrm{kpc}}
\providecommand{\lP}{}\renewcommand{\lP}{\ell_{P}}
\providecommand{\mD}{}\renewcommand{\mD}{\mathbf{m}_{D}}
\providecommand{\mb}{}\renewcommand{\mb}{\bar m}
\providecommand{\mnu}{}\renewcommand{\mnu}{\Sigma m_{\nu}}
\providecommand{\nInf}{}\renewcommand{\nInf}{n_\infty}
\providecommand{\nO}{}\renewcommand{\nO}{n_0}
\providecommand{\ninf}{}\renewcommand{\ninf}{n_{\infty}}
\providecommand{\nz}{}\renewcommand{\nz}{n_{0}}
\providecommand{\proj}{}\renewcommand{\proj}{\operatorname{proj}}
\providecommand{\pv}{}\renewcommand{\pv}{\pi}
\providecommand{\rad}{}\renewcommand{\rad}{\mathrm{rad}}
\providecommand{\sgn}{}\renewcommand{\sgn}{\operatorname{sgn}}
\providecommand{\shA}{}\renewcommand{\shA}{\mathcal{A}}
\providecommand{\shC}{}\renewcommand{\shC}{\mathcal{C}}
\providecommand{\shE}{}\renewcommand{\shE}{\mathcal{E}}
\providecommand{\shJ}{}\renewcommand{\shJ}{\mathcal{J}}
\providecommand{\shK}{}\renewcommand{\shK}{\mathcal{K}}
\providecommand{\shL}{}\renewcommand{\shL}{L}
\providecommand{\shW}{}\renewcommand{\shW}{\mathcal{W}}
\providecommand{\sll}{}\renewcommand{\sll}{\mathfrak{sl}}
\providecommand{\so}{}\renewcommand{\so}{\mathfrak{so}}
\providecommand{\su}{}\renewcommand{\su}{\mathfrak{su}}
\providecommand{\supp}{}\renewcommand{\supp}{\operatorname{supp}}
\providecommand{\tr}{}\renewcommand{\tr}{\operatorname{tr}}
\providecommand{\uu}{}\renewcommand{\uu}{\mathfrak{u}}
\providecommand{\br}{}\renewcommand{\br}[1]{\langle #1\rangle}
\providecommand{\bra}{}\renewcommand{\bra}[1]{\langle #1|}
\providecommand{\file}{}\renewcommand{\file}[1]{\texttt{#1}}
\providecommand{\hp}{}\renewcommand{\hp}[2]{( #1 \,|\, #2 )}
\providecommand{\ip}{}\renewcommand{\ip}[2]{\langle #1, #2\rangle}
\providecommand{\ket}{}\renewcommand{\ket}[1]{|#1\rangle}
\providecommand{\rel}{}\renewcommand{\rel}[2]{D\bigl(#1 \,\|\, #2\bigr)}
\providecommand{\rev}{}\renewcommand{\rev}[1]{\widetilde{#1}}
\providecommand{\slashed}{}\renewcommand{\slashed}[1]{{#1}\mkern-10mu/}
\title{\bfseries Zero divisors and obstructions}
\author{Leon Butler\\ \normalsize\href{mailto:lbutler.math@gmail.com}{\texttt{lbutler.math@gmail.com}}\\ \normalsize ORCID: \href{https://orcid.org/0009-0006-4449-5617}{0009-0006-4449-5617}}
\date{Generated by the Transport Lab, 2 November 2026}
\begin{document}
\maketitle
\begin{abstract}
\normalsize
Division is the operation this whole series turns on, and this trail follows where it breaks down.

A null vector, a nonzero vector $v$ whose square $v^{2}$ is zero, has no inverse, even inside a larger algebra. With a nondegenerate form, every failure of division is one of three kinds: a null vector, an algebra split into two halves, or a projection.

The projection kind repairs itself. The null kind cannot be repaired element by element while respecting the symmetries, however the algebra is enlarged. And every repair creates failures of its own.
\end{abstract}
\newpage
\noindent \textbf{Authorship and disclosure. }The founding idea of the series is the author's: that division by zero in a Clifford algebra should be read not as a prohibition but as a transport into a higher-dimensional algebra, and that the transported division need not be accompanied by an inverse operation. The direction of investigation and the critical review of the results are likewise the author's. The mathematical development and the drafting of the text were carried out in extended collaboration with Claude Fable 5 (Anthropic), an AI system.

\noindent\fbox{\parbox{\dimexpr\linewidth-2\fboxsep-2\fboxrule}{\textbf{This is not a deposited paper.} It was generated by the Transport Lab on 2 November 2026 and assembled from approved parts of its store. It follows results across several papers of the series; each is stated with its proof and credited to the paper that first proved it. The official versions of those papers are on Zenodo, in the Transport Papers community, \url{https://zenodo.org/communities/transport-papers}.}}
\newpage

\section*{Introduction}
The walk starts with the simplest element that cannot be divided by, maps where each kind of failure occurs, shows the list of kinds is complete, and ends with what can and cannot be repaired.

\section{Zero divisors and obstructions}

These results ask when and why division fails in a Clifford algebra.

A null vector is a nonzero $v$ with $v^{2} = 0$.

If $v$ had an inverse $u$, then $v = v(vu) = v^{2} u = 0$, so no inverse exists on either side, in $\Cl(p,q)$ or in any larger associative algebra containing it. This is the first and simplest failure of division, and the rest of the walk grows from it.

\begin{lemma}[A null vector has no inverse]\label{c:transport.noinverse}
A null vector is a nonzero direction whose square is zero. Such an element can never be divided by: if it had an inverse on either side, multiplying through would force the vector itself to be zero. The argument uses only the fact that its square is zero, so it still applies in any larger algebra of the same kind that contains the original one.
\emph{First proved in} T1, Lem 2.1 (\href{https://doi.org/10.5281/zenodo.21743863}{doi:10.5281/zenodo.21743863}).
\end{lemma}
\noindent\emph{Rests on:} \ref{c:std-algebra.vector-squares}.

Next comes the question of where such vectors exist.

They exist exactly when the signature is mixed, \[ pq \neq 0 \] that is, when there are both positive and negative directions, as in spacetime. They then form a cone of dimension $n - 1$, like a light cone, and every one of them is a two-sided zero divisor.

\begin{proposition}[Existence locus of type N]\label{c:transport.existence-locus-type-n}
In a space with both positive and negative directions, like the space-time of relativity, there are nonzero directions of zero length, called null directions. This result says they exist exactly when both kinds of direction are present, and that they then form a cone, much like a light cone. Each such direction multiplied by itself gives zero, so it can never be divided by, even in a larger algebra that contains the original one.
\emph{First proved in} T1, Prop 3.3 (\href{https://doi.org/10.5281/zenodo.21743863}{doi:10.5281/zenodo.21743863}).
\end{proposition}
\noindent\emph{Rests on:} \ref{c:transport.noinverse}, \ref{c:std-algebra.orthogonal-anticommute}.

With one failure in view, it is natural to ask which algebras escape failure altogether.

Only three do: \[ \Cl(0,0) = \R, \quad \Cl(0,1) = \C, \quad \Cl(0,2) = \HH \] the reals, the complex numbers and the quaternions. Every other real Clifford algebra has nonzero elements whose product is zero.

\begin{proposition}[The exceptional division algebras]\label{c:transport.exceptional-division-algebras}
In most Clifford algebras some nonzero elements multiply to zero, so division sometimes fails. This result lists the only exceptions: the real numbers, the complex numbers and the quaternions. Allowing pairs of these as well gives exactly the algebras in which no nonzero element has a power equal to zero.
\emph{First proved in} T1, Prop 3.2 (\href{https://doi.org/10.5281/zenodo.21743863}{doi:10.5281/zenodo.21743863}).
\end{proposition}
\noindent\emph{Rests on:} \ref{c:transport.classification-real-clifford-algebras}, \ref{c:std-algebra.blade-basis}.

A second way for division to fail is for the whole algebra to come apart into two independent halves.

This happens exactly when $p - q \equiv 1 \pmod 4$, and the halves are picked out by the two central idempotents \[ \varepsilon\_\pm = \tfrac12\,(1 \pm \omega) \] built from the pseudoscalar $\omega$, the product of all the basis vectors. They satisfy $\varepsilon_+ \varepsilon_- = 0$ and $\varepsilon_+ + \varepsilon_- = 1$, and the grade involution swaps them, but no inner automorphism can.

\begin{theorem}[Existence and classification of type C]\label{c:transport.existence-classification-type-c}
Some Clifford algebras secretly come in two independent halves, and the elements that pick out each half are called central idempotents. This result pins down exactly when that happens: precisely when the number of positive directions minus the number of negative ones leaves remainder 1 on division by four. In that case there are exactly two such elements, built from the product of all the basis directions. Reversing every direction swaps the two halves, but no conjugation by an element of the algebra itself can do so, since these elements commute with everything.
\emph{First proved in} T1, Thm 3.6 (\href{https://doi.org/10.5281/zenodo.21743863}{doi:10.5281/zenodo.21743863}).
\end{theorem}
\noindent\emph{Rests on:} \ref{c:transport.centre}, \ref{c:transport.classification-real-clifford-algebras}, \ref{c:transport.sign}, \ref{c:std-algebra.involution-properties}.

The third kind, type P, is an idempotent that does not commute with everything, like a projection onto part of a space.

Such elements are almost everywhere: for an orthogonal pair $a, b$ with $a^{2} b^{2} = -1$, the bivector $u = ab$ squares to one and gives \[ \tfrac12(1 + u) \] a non-central idempotent. Only five small algebras, built from the reals, the complex numbers and the quaternions, lack them.

\begin{proposition}[Existence locus of type P]\label{c:transport.existence-locus-type-p}
An element that equals its own square but does not commute with everything is a type-P obstruction. Such elements exist in every real Clifford algebra except five small ones built from the real numbers, the complex numbers and the quaternions. A simple recipe produces one from a space-like and a time-like direction at right angles.
\emph{First proved in} T1, Prop 3.8 (\href{https://doi.org/10.5281/zenodo.21743863}{doi:10.5281/zenodo.21743863}).
\end{proposition}
\noindent\emph{Rests on:} \ref{c:transport.classification-real-clifford-algebras}, \ref{c:transport.exceptional-division-algebras}, \ref{c:std-algebra.orthogonal-anticommute}.

With three kinds named, the question is whether there are any others.

The key step is that each element $a$ of a nondegenerate Clifford algebra has a partial inverse $x$ with $axa = a$. For a zero divisor this produces idempotents $e = xa$ and $f = ax$ with \[ a(1 - e) = 0, \qquad (1 - f)a = 0 \] which pin down where the failure sits.

\begin{lemma}[Regularity localises zero divisors at idempotents]\label{c:transport.regularity-localises-zero-divisors-idempotents}
In the algebras of the series every element has a weak kind of inverse that undoes it partway. Applying that weak inverse to an element that cannot be divided by produces elements equal to their own squares, which pin down exactly where the failure sits. This is the step that turns every failure of division into a statement about such elements.
\emph{First proved in} T1, Lem 2.8 (\href{https://doi.org/10.5281/zenodo.21743863}{doi:10.5281/zenodo.21743863}).
\end{lemma}
\noindent\emph{Rests on:} \ref{c:transport.classification-real-clifford-algebras}.

That step turns every failure of division into a statement about idempotents.

From there, whenever the form is nondegenerate, every zero divisor traces back to one of the three kinds already named, so the list is complete. The central kind can occur only when $p - q \equiv 1 \pmod 4$; otherwise every annihilating idempotent is non-central.

\begin{theorem}[Exhaustiveness: every zero divisor is annihilated by an idempotent]\label{c:transport.exhaustiveness-every-zero-divisor-annihilated-idempotent}
Whenever division fails in a Clifford algebra with a nondegenerate notion of length, the failure can be traced to an element that equals its own square. This shows that the three kinds of failure already named cover every case. Which kinds occur depends only on the numbers of space and time directions.
\emph{First proved in} T1, Thm 3.9 (\href{https://doi.org/10.5281/zenodo.21743863}{doi:10.5281/zenodo.21743863}).
\end{theorem}
\noindent\emph{Rests on:} \ref{c:transport.classification-real-clifford-algebras}, \ref{c:transport.existence-classification-type-c}, \ref{c:transport.existence-locus-type-n}, \ref{c:transport.regularity-localises-zero-divisors-idempotents}.

Once the kinds are known, one can ask which of them can be repaired.

The projection kind repairs itself: each idempotent is its own partial inverse, \[ e\,e\,e = e \] so the assignment $\sigma(e) = e$ satisfies the required identities and respects every symmetry. Nothing needs to be added to the algebra.

\begin{proposition}[Type P is self-repairing]\label{c:transport.type-p-self-repairing}
Some zero divisors are projection-like elements, ones that equal their own square. For these, the repair being sought is already available: each such element can serve as its own partner and satisfies the required identities. Because the rule is simply to take the element itself, it automatically respects every symmetry, so nothing needs to be added to the algebra.
\emph{First proved in} T1, Prop 6.1 (\href{https://doi.org/10.5281/zenodo.21743863}{doi:10.5281/zenodo.21743863}).
\end{proposition}

The null kind is the opposite case.

A null vector $v$ has many partial inverses $\sigma(v)$ with \[ v\,\sigma(v)\,v = v \] but in mixed signature with $n \geq 3$ there is no way to choose one for each null $v$ that is equivariant, even under $\Spin^0(p,q)$. Enlarging the algebra by any number of new directions does not help.

Any repair of this kind must therefore work differently, and the completion trail takes that up.

\begin{theorem}[No-go at the level of elements]\label{c:transport.no-go-level-elements}
A light-like direction has no inverse, but it has many partial inverses. This theorem shows that there is no way to choose one for every light-like direction that respects the symmetries of the geometry, in any space with both space and time directions and at least three dimensions. Enlarging the algebra by any number of new directions does not help.
\emph{First proved in} T1, Thm 4.2 (\href{https://doi.org/10.5281/zenodo.21743863}{doi:10.5281/zenodo.21743863}).
\end{theorem}
\noindent\emph{Rests on:} \ref{c:transport.splitting-along-hyperbolic-plane}, \ref{c:primer.commutation-rule}, \ref{c:std-algebra.witt-decomposition}.

A last question is whether repairing ever finishes.

It does not: a completion by the first rule has mixed signature and so new null vectors, and one by the second is a matrix algebra $\Mat_2(\Cl^0(p,q))$ with new projections. So the process never ends in an algebra free of zero divisors.

\begin{proposition}[No transport terminates]\label{c:transport.no-transport-terminates}
Each repair in the transport programme enlarges the algebra to fix a failure of division. This result shows that the enlarged algebra always has failures of its own: one repair adds new light-like directions, the other produces projections and elements that square to zero. So the process never ends in an algebra where every nonzero element can be divided by.
\emph{First proved in} T1, Prop 7.8 (\href{https://doi.org/10.5281/zenodo.21743863}{doi:10.5281/zenodo.21743863}).
\end{proposition}
\noindent\emph{Rests on:} \ref{c:transport.existence-locus-type-n}, \ref{c:transport.matrix-units-peirce-isomorphism}, \ref{c:transport.rule-2-ideal-fusion-sharp-form}.

\section{Results used from elsewhere}
The proofs above use the following results of the series, stated briefly; each is proved in the paper named.

\begin{fact}[The commutation rule]\label{c:primer.commutation-rule}
If $x$ lies in the subalgebra generated by vectors orthogonal to $v$, then $v x = \hat{x} v$ and $x v = v \hat{x}$.
\emph{Source:} T0, Lem 9.1 (\href{https://doi.org/10.5281/zenodo.22714501}{doi:10.5281/zenodo.22714501}).
\end{fact}

\begin{fact}[The blade basis, degenerate forms included]\label{c:std-algebra.blade-basis}
For any basis of $V$ the $2^{n}$ increasing products $e_I$ form a basis of $\Cl(V, Q)$, degenerate forms included; so $V \hookrightarrow \Cl(V, Q)$.
\emph{Source:} the algebra core of the Transport Lab, proved there with an exact check; standard, see C. Chevalley, The Algebraic Theory of Spinors, Columbia University Press (1954), Ch. II (the Clifford algebra of a quadratic space).
\end{fact}

\begin{fact}[The three involutions on each grade]\label{c:std-algebra.involution-properties}
On $\Cl^{k}$: $\hat{x} = (-1)^{k}x$, $\tilde{x} = (-1)^{k(k-1)/2}x$, $\bar{x} = (-1)^{k(k+1)/2}x$; the first is an automorphism, the others anti-automorphisms.
\emph{Source:} the algebra core of the Transport Lab, proved there with an exact check; standard, see P. Lounesto, Clifford Algebras and Spinors, 2nd ed., London Mathematical Society Lecture Note Series 286, Cambridge University Press (2001), Ch. 14 (definitions of the Clifford algebra).
\end{fact}

\begin{fact}[Orthogonal vectors anticommute]\label{c:std-algebra.orthogonal-anticommute}
$uv = -vu$ iff $B(u,v) = 0$; in an orthogonal basis $e_I e_J = \pm \prod_{I \cap J} Q(e_i)\; e_{I \triangle J}$ with the sorting sign.
\emph{Source:} the algebra core of the Transport Lab, proved there with an exact check; standard, see C. Doran and A. Lasenby, Geometric Algebra for Physicists, Cambridge University Press (2003), Ch. 4 (foundations of geometric algebra).
\end{fact}

\begin{fact}[The product of two vectors]\label{c:std-algebra.vector-squares}
$v^{2} = Q(v)$ and $uv + vu = 2B(u, v)$, so $uv = B(u, v) + \tfrac12(uv - vu)$.
\emph{Source:} the algebra core of the Transport Lab, proved there with an exact check; standard, see C. Doran and A. Lasenby, Geometric Algebra for Physicists, Cambridge University Press (2003), Ch. 4 (foundations of geometric algebra).
\end{fact}

\begin{fact}[Hyperbolic pairs and the Witt decomposition]\label{c:std-algebra.witt-decomposition}
A null vector in a nondegenerate plane has a unique hyperbolic partner, and $V = H_1 \perp \dots \perp H_{\min(p,q)} \perp A \perp \operatorname{rad} V$ with $A$ definite; maximal totally null subspaces have dimension $\min(p,q) + r$.
\emph{Source:} the algebra core of the Transport Lab, proved there with an exact check; standard, see T. Y. Lam, Introduction to Quadratic Forms over Fields, Graduate Studies in Mathematics 67, American Mathematical Society (2005), Ch. I (hyperbolic planes, the Witt decomposition, Witt's theorems).
\end{fact}

\begin{fact}[The centre]\label{c:transport.centre}
For $n = p + q \geq 1$ the centre of $\Cl(p,q)$ is $\R \cdot 1$ when $n$ is even and $\R \cdot 1 \oplus \R \cdot \omega$ when $n$ is odd, with $\omega$ the pseudoscalar.
\emph{Source:} T1, Lem 2.6 (\href{https://doi.org/10.5281/zenodo.21743863}{doi:10.5281/zenodo.21743863}).
\end{fact}

\begin{fact}[Classification of the real Clifford algebras]\label{c:transport.classification-real-clifford-algebras}
The real Clifford algebra $\Cl(p,q)$ is a matrix algebra over $\R$, $\C$ or $\HH$, or a sum of two such, according to $p - q \bmod 8$, and it is simple except when $p - q \equiv 1 \pmod 4$.
\emph{Source:} M. F. Atiyah, R. Bott and A. Shapiro, Clifford modules (1964); I. R. Porteous, Clifford Algebras and the Classical Groups (1995); P. Lounesto, Clifford Algebras and Spinors (2001).
\end{fact}

\begin{fact}[Matrix units, and the Peirce isomorphism]\label{c:transport.matrix-units-peirce-isomorphism}
The projectors $\varepsilon_\pm$ and the adjoined generator $\nu$ form a full system of $2 \times 2$ matrix units, so $B_\gamma \cong \Mat_2(E_{11} B_\gamma E_{11})$.
\emph{Source:} T1, Lem 5.4 (\href{https://doi.org/10.5281/zenodo.21743863}{doi:10.5281/zenodo.21743863}).
\end{fact}

\begin{fact}[Rule 2: ideal fusion, sharp form]\label{c:transport.rule-2-ideal-fusion-sharp-form}
For $p - q \equiv 1 \pmod 4$, adjoining one generator $\nu$ with $Q(\nu) = \pm 1$ gives the simple algebra $\Mat_2(\Cl^0(p,q))$ with centre $\R$, the same for either sign.
\emph{Source:} T1, Thm 5.5 (\href{https://doi.org/10.5281/zenodo.21743863}{doi:10.5281/zenodo.21743863}).
\end{fact}

\begin{fact}[The sign of $\omega^2$]\label{c:transport.sign}
The pseudoscalar of $\Cl(p,q)$ satisfies $\omega^2 = (-1)^{(p-q)(p-q-1)/2}$, so $\omega^2 = +1$ exactly when $p - q \equiv 0$ or $1 \pmod 4$.
\emph{Source:} T1, Lem 2.7 (\href{https://doi.org/10.5281/zenodo.21743863}{doi:10.5281/zenodo.21743863}).
\end{fact}

\begin{fact}[Splitting along a hyperbolic plane]\label{c:transport.splitting-along-hyperbolic-plane}
Along a hyperbolic plane spanned by $(v, w_0)$, every element of $\Cl(H_0 \perp U')$ expands uniquely as $\alpha + v\beta + w_0\gamma + v w_0\delta$ with coefficients in $\Cl(U')$, and left multiplication by $v$ is injective on $\Cl(U')$.
\emph{Source:} T1, Lem 2.4 (\href{https://doi.org/10.5281/zenodo.21743863}{doi:10.5281/zenodo.21743863}).
\end{fact}

\section*{Conclusion}
Division fails in exactly three ways in a Clifford algebra with a nondegenerate form, and the list is complete.

No repair ends in an algebra free of failures, so obstructions are something to move and manage rather than abolish. This trail rests on no postulate: each step is proved here or cited from results proved elsewhere, including the classical classification of real Clifford algebras.

\end{document}